\begin{theorem}[Constant Multiple Rule]
\label{thm:constant-multiple-rule}
Let $A \subseteq \mathbb{R}$, let $f : A \to \mathbb{R}$, let $c \in A$ be a limit point of $A$, and let $\alpha \in \mathbb{R}$. If $f$ is differentiable at $c$, then $\alpha f$ is differentiable at $c$, and
\[
(\alpha f)'(c)=\alpha f'(c).
\]
\end{theorem}

\begin{theorem}[Sum Rule]
\label{thm:sum-rule}
Let $A \subseteq \mathbb{R}$, let $f,g : A \to \mathbb{R}$, and let $c \in A$ be a limit point of $A$. If $f$ and $g$ are differentiable at $c$, then $f+g$ is differentiable at $c$, and
\[
(f+g)'(c)=f'(c)+g'(c).
\]
\end{theorem}

\begin{theorem}[Product Rule]
\label{thm:product-rule}
Let $A \subseteq \mathbb{R}$, let $f,g : A \to \mathbb{R}$, and let $c \in A$ be a limit point of $A$. If $f$ and $g$ are differentiable at $c$, then $fg$ is differentiable at $c$, and
\[
(fg)'(c)=f'(c)g(c)+f(c)g'(c).
\]
\end{theorem}

\begin{theorem}[Quotient Rule]
\label{thm:quotient-rule}
Let $A \subseteq \mathbb{R}$, let $f,g : A \to \mathbb{R}$, and let $c \in A$ be a limit point of $A$. Suppose that $f$ and $g$ are differentiable at $c$ and that $g(c)\neq 0$. Then $f/g$ is differentiable at $c$, and
\[
\left(\frac{f}{g}\right)'(c)=\frac{f'(c)g(c)-f(c)g'(c)}{(g(c))^2}.
\]
\end{theorem}

\begin{corollary}[Finite Sum Rule]
\label{cor:finite-sum-rule}
Let $A \subseteq \mathbb{R}$, let $f_1,\dots,f_n : A \to \mathbb{R}$, let $\alpha_1,\dots,\alpha_n \in \mathbb{R}$, and let $c \in A$ be a limit point of $A$. If each $f_i$ is differentiable at $c$, then $\sum_{i=1}^n \alpha_i f_i$ is differentiable at $c$, and
\[
\left(\sum_{i=1}^n \alpha_i f_i\right)'(c)=\sum_{i=1}^n \alpha_i f_i'(c).
\]
\end{corollary}

\begin{corollary}[Extended Product Rule]
\label{cor:extended-product-rule}
Let $A \subseteq \mathbb{R}$, let $f_1,\dots,f_n : A \to \mathbb{R}$, and let $c \in A$ be a limit point of $A$. If each $f_i$ is differentiable at $c$, then the product $\prod_{i=1}^n f_i$ is differentiable at $c$, and
\[
\left(\prod_{i=1}^n f_i\right)'(c)
=
\sum_{k=1}^n
\left(
f_1(c)\cdots f_{k-1}(c)\,f_k'(c)\,f_{k+1}(c)\cdots f_n(c)
\right).
\]
\end{corollary}

\begin{corollary}[Power Rule for Integer Powers of a Function]
\label{cor:power-rule-special-case}
Let $A \subseteq \mathbb{R}$, let $f : A \to \mathbb{R}$, let $c \in A$ be a limit point of $A$, and let $n \in \mathbb{N}$. If $f$ is differentiable at $c$, then $f^n$ is differentiable at $c$, and
\[
(f^n)'(c)=n(f(c))^{n-1}f'(c).
\]
\end{corollary}

\begin{theorem}[Finite Linear Combination Rule]
\label{thm:finite-linear-combination-rule}
Let $A \subseteq \mathbb{R}$, let $f_1,\dots,f_n : A \to \mathbb{R}$, let $\alpha_1,\dots,\alpha_n \in \mathbb{R}$, and let $c \in A$ be a limit point of $A$. If each $f_i$ is differentiable at $c$, then the function
\[
x \mapsto \sum_{i=1}^n \alpha_i f_i(x)
\]
is differentiable at $c$, and
\[
\left(\sum_{i=1}^n \alpha_i f_i\right)'(c)=\sum_{i=1}^n \alpha_i f_i'(c).
\]
\end{theorem}



\begin{theorem}[Constant Multiple Rule on an Interval]
\label{thm:constant-multiple-rule}
Let $I \subseteq \mathbb{R}$ be an interval, let $f : I \to \mathbb{R}$, and let $\alpha \in \mathbb{R}$. If $f$ is differentiable on $I$, then $\alpha f$ is differentiable on $I$, and
\[
(\alpha f)'(x)=\alpha f'(x)
\qquad \text{for all } x \in I.
\]
\end{theorem}

\begin{theorem}[Sum Rule on an Interval]
\label{thm:sum-rule}
Let $I \subseteq \mathbb{R}$ be an interval, and let $f,g : I \to \mathbb{R}$. If $f$ and $g$ are differentiable on $I$, then $f+g$ is differentiable on $I$, and
\[
(f+g)'(x)=f'(x)+g'(x)
\qquad \text{for all } x \in I.
\]
\end{theorem}

\begin{theorem}[Product Rule on an Interval]
\label{thm:product-rule}
Let $I \subseteq \mathbb{R}$ be an interval, and let $f,g : I \to \mathbb{R}$. If $f$ and $g$ are differentiable on $I$, then $fg$ is differentiable on $I$, and
\[
(fg)'(x)=f'(x)g(x)+f(x)g'(x)
\qquad \text{for all } x \in I.
\]
\end{theorem}

\begin{theorem}[Quotient Rule on an Interval]
\label{thm:quotient-rule}
Let $I \subseteq \mathbb{R}$ be an interval, and let $f,g : I \to \mathbb{R}$. Suppose that $f$ and $g$ are differentiable on $I$ and that $g(x)\neq 0$ for all $x \in I$. Then $f/g$ is differentiable on $I$, and
\[
\left(\frac{f}{g}\right)'(x)=\frac{f'(x)g(x)-f(x)g'(x)}{(g(x))^2}
\qquad \text{for all } x \in I.
\]
\end{theorem}

\begin{corollary}[Finite Sum Rule on an Interval]
\label{cor:finite-sum-rule}
Let $I \subseteq \mathbb{R}$ be an interval, let $f_1,\dots,f_n : I \to \mathbb{R}$, and let $\alpha_1,\dots,\alpha_n \in \mathbb{R}$. If each $f_i$ is differentiable on $I$, then $\sum_{i=1}^n \alpha_i f_i$ is differentiable on $I$, and
\[
\left(\sum_{i=1}^n \alpha_i f_i\right)'(x)=\sum_{i=1}^n \alpha_i f_i'(x)
\qquad \text{for all } x \in I.
\]
\end{corollary}

\begin{corollary}[Extended Product Rule on an Interval]
\label{cor:extended-product-rule}
Let $I \subseteq \mathbb{R}$ be an interval, and let $f_1,\dots,f_n : I \to \mathbb{R}$. If each $f_i$ is differentiable on $I$, then $\prod_{i=1}^n f_i$ is differentiable on $I$, and
\[
\left(\prod_{i=1}^n f_i\right)'(x)
=
\sum_{k=1}^n
\left(
f_1(x)\cdots f_{k-1}(x)\,f_k'(x)\,f_{k+1}(x)\cdots f_n(x)
\right)
\qquad \text{for all } x \in I.
\]
\end{corollary}

\begin{corollary}[Power Rule for Integer Powers of a Function on an Interval]
\label{cor:power-rule-special-case}
Let $I \subseteq \mathbb{R}$ be an interval, let $f : I \to \mathbb{R}$, and let $n \in \mathbb{N}$. If $f$ is differentiable on $I$, then $f^n$ is differentiable on $I$, and
\[
(f^n)'(x)=n(f(x))^{n-1}f'(x)
\qquad \text{for all } x \in I.
\]
\end{corollary}

\begin{theorem}[Finite Linear Combination Rule on an Interval]
\label{thm:finite-linear-combination-rule}
Let $I \subseteq \mathbb{R}$ be an interval, let $f_1,\dots,f_n : I \to \mathbb{R}$, and let $\alpha_1,\dots,\alpha_n \in \mathbb{R}$. If each $f_i$ is differentiable on $I$, then the function
\[
x \mapsto \sum_{i=1}^n \alpha_i f_i(x)
\]
is differentiable on $I$, and
\[
\left(\sum_{i=1}^n \alpha_i f_i\right)'(x)=\sum_{i=1}^n \alpha_i f_i'(x)
\qquad \text{for all } x \in I.
\]
\end{theorem}



\begin{remark*}[Global differentiability]
The pointwise statements above — each asserting differentiability and a
derivative formula at a single point $c \in I$ — extend immediately to global
statements when the hypotheses hold at every point. If $f$ is differentiable
on $I$ (meaning at every $c \in I$), then $f'$ is itself a function
$f' : I \to \mathbb{R}$, and differentiation acts as an operator on the space
of differentiable functions. The algebra rules then read:
\[
  (\lambda f)' = \lambda f', \qquad
  (f \pm g)' = f' \pm g', \qquad
  (fg)' = f'g + fg', \qquad
  \left(\frac{f}{g}\right)' = \frac{f'g - fg'}{g^2},
\]
where equality holds as functions on $I$, and $g(x) \neq 0$ for all $x \in I$
in the quotient case. These are not separate theorems — they are the pointwise
results collected into function notation. The Finite Linear Combination Rule
in this global form is the statement that differentiation is a linear operator
on the vector space of differentiable functions on $I$.
\end{remark*}


\begin{remark*}[Global form]
When $f$ is strictly monotone on $I$ (continuity of $f$ is not separately
required — it follows from strict monotonicity on an interval up to endpoint
behaviour) and differentiable on $I$ with $f'(x) \neq 0$ for all $x \in I$,
the pointwise result holds at every $c \in I$ simultaneously. The inverse
$g : J \to I$ is then differentiable on $J$ and
\[
  g' = \frac{1}{f' \circ g},
\]
where equality holds as functions on $J$. This is the global statement obtained
by applying the pointwise theorem at each $d \in J$ with $c = g(d)$.
\end{remark*}

